% s-351-p351-m002; original PDF 351
\Aphorism{s-351-p351-m002}{402}
Ich sah Ameisen auf pfadlosem Waldboden, den Jahre lang kein Mensch betrat. Die „Weltanschauung“ dieser Thiere mag in der Anpassung an die Umwelt noch so fein, in der Sammlung älterer Erfahrungen noch so umfassend sein, sie mag den Besonderheiten der näheren und ferneren Umgebung, den klimatischen Schwankungen in Tag und Jahr, allen normalen Vorfällen des einzelnen und socialen Ameisenlebens gerecht werden, — die Erscheinung des Menschen, so selten, so unberechenbar, wird sicherlich aus diesem Rahmen heraustreten und den Eindruck einer Discontinuität, einer irrationalen, übernatürlichen Durchbrechung aller Gesetze hinterlassen. — Was aber hindert uns, den Menschen und seine Welt auch in solch einem einsamen Winkel des Alls gelegen anzunehmen, den aller Jahrmyriaden irgend ein gewaltsam tappender Tölpel von Gott betritt? Mit andern Worten: ist unsere „Weltanschauung“, die nicht an plötzliche Eingriffe in den gesetzlichen Gang der Begebenheiten glauben will, nicht vielleicht eine Local- und Intermezzo-Philosophie, gewachsen in der ruhigen Zwischenzeit zwischen zwei solchen Eingriffen? Also nur ein Bewusstsein von Gestern und Heute, vielleicht schon nicht mehr von Morgen, die wissenschaftlich ausgedrückte Arglosigkeit und Vertrauensseligkeit unseres Ameisenhaufens, der sich vom letzten Fusstritt bis zum Vergessen erholt hat und nicht ahnt, wie fern oder nah ihm der nächste ist?
\par

